% ============================================================================
% 부록
% ============================================================================

\chapter{재현 절차}\label{app:repro}

본 연구의 모든 코드와 수치 결과는 연구 저장소에 있다. 원본 데이터셋 영상과 모델 가중치는 용량과 재배포
문제로 저장소에 포함하지 않았다. 아래 명령은 저장소 최상위에서 실행한다.

\section{환경 구성}

\begin{small}
\begin{verbatim}
# 1) ultralytics 기반 커밋에 본 연구의 변경분(패치)을 적용
git clone https://github.com/ultralytics/ultralytics.git ultralytics_thyroid
cd ultralytics_thyroid
git checkout 4dcc98a6be09d058c206f2756f3ee2966d1c4b71
git apply ../ultralytics_patches/thyroid.patch
pip install -e . && cd ..
# 2) 부가 의존성
pip install onnx onnxruntime-gpu openvino nncf scipy
\end{verbatim}
\end{small}

패치는 \BRM{} 헤드(\texttt{nn/modules/head.py}), 경계 인식 손실(\texttt{utils/loss.py}), 모델 구성 파일
(\texttt{cfg/models/11/yolo11-seg-brm.yaml})과 이들의 등록 코드를 담고 있다. 경계 손실 계수는 환경 변수
\texttt{BRM\_HEAD\_GAIN}($g$, 기본 1.0)과 \texttt{BRM\_LAMBDA}($\lambda$, 기본 2.0)로 바꿀 수 있다.

\section{실행 순서}

\begin{small}
\begin{verbatim}
python scripts/convert_to_yoloseg.py        # 마스크 -> YOLO-seg 다각형 (3.4절)
python scripts/make_splits.py               # 5-fold / LODO / held-out 분할 (3.5절)
bash   scripts/run_baseline_rest.sh         # 베이스라인 nano·small 학습 (5.4절)
bash   scripts/run_phase3_sweep.sh          # 경계 손실 계수 스윕 (5.5절)
python scripts/sweep_eval.py                # 스윕 결과표, held-out 기준 (5.5절)
python scripts/sweep_eval.py --split val    # 같은 모델의 검증 분할 기준 사후 점검
bash   scripts/run_phase3_final.sh          # 제안 모델 5-fold + LODO 학습 (5.6절)
bash   scripts/run_lodo_seeds.sh            # LODO 시드 43·44 반복 (5.9절)
python scripts/stats_test.py --part all     # 대응표본 검정 (5.8절, 5.9절)
python scripts/export_quant.py --fold N     # 정밀도별 변환, N = 0..4 (4.6절)
python scripts/eval_quant.py --folds 0,1,2,3,4   # 정밀도별 평가 (5.10절)
python scripts/final_table.py               # 단계별 누적 결과 (5.11절)
python scripts/make_thesis.py --zip         # 본 논문의 표·그림 재생성
\end{verbatim}
\end{small}

학습 스크립트는 같은 이름의 출력 폴더를 덮어쓰며 이어 쓰도록 설정되어 있어, 같은 실행을 두 번 띄우면
두 학습의 기록과 가중치가 섞인다. 실제로 이 문제로 제안 모델의 fold 0--2 결과가 오염된 적이 있어, 해당 fold를
처음부터 다시 학습하였고 재실행 스크립트에 기존 폴더가 있으면 실행을 거부하는 장치를 넣었다
(\appref{app:impl}). 평가는 반드시 학습 드라이버의 종료 기록을 확인한 뒤 수행해야 한다.

\chapter{상세 결과}\label{app:detail}

\section{fold별 결과}
% 자동 생성 파일 — scripts/make_thesis.py 가 덮어쓴다. 직접 고치지 말 것.
\begin{table}[htbp]
\centering
\footnotesize
\caption[fold별 상세 결과 (held-out TN3K test).]{fold별 상세 결과 (held-out TN3K test). 기준 = 베이스라인 nano, 제안 = BRM nano.}
\label{tab:perfold}
\begin{adjustbox}{max width=\linewidth}
\begin{tabular}{lcccccccc}
\toprule
\multirow{2}{*}{fold} & \multicolumn{2}{c}{Mask mAP$_{50:95}$} & \multicolumn{2}{c}{Dice} & \multicolumn{2}{c}{HD95 (\%)} & \multicolumn{2}{c}{BIoU} \\
\cmidrule(lr){2-3}\cmidrule(lr){4-5}\cmidrule(lr){6-7}\cmidrule(lr){8-9}
 & 기준 & 제안 & 기준 & 제안 & 기준 & 제안 & 기준 & 제안 \\
\midrule
fold 0 & 0.6215 & 0.6271 & 0.8095 & 0.8344 & 10.656 & 7.034 & 0.4265 & 0.4374 \\
fold 1 & 0.6256 & 0.6186 & 0.8206 & 0.8309 & 9.447 & 8.829 & 0.4320 & 0.4422 \\
fold 2 & 0.6398 & 0.6266 & 0.8170 & 0.8301 & 11.015 & 8.429 & 0.4448 & 0.4381 \\
fold 3 & 0.6219 & 0.6215 & 0.8221 & 0.8237 & 8.578 & 8.932 & 0.4351 & 0.4395 \\
fold 4 & 0.6272 & 0.6398 & 0.8284 & 0.8306 & 8.444 & 8.669 & 0.4408 & 0.4496 \\
\midrule
평균 & 0.6272 & 0.6267 & 0.8195 & 0.8300 & 9.628 & 8.378 & 0.4358 & 0.4414 \\
\bottomrule
\end{tabular}
\end{adjustbox}
\end{table}

% 자동 생성 파일 — scripts/make_thesis.py 가 덮어쓴다. 직접 고치지 말 것.
\begin{table}[htbp]
\centering
\small
\caption[fold별 학습 epoch (총 학습 epoch / 체크포인트 epoch).]{fold별 학습 epoch (총 학습 epoch / 체크포인트 epoch). 최대 100 epoch, 검증 fitness가 20 epoch 동안 개선되지 않으면 조기 종료하였으므로, 총 epoch이 100보다 작은 경우 체크포인트 epoch보다 정확히 20 크다.}
\label{tab:epochs}
\begin{adjustbox}{max width=\linewidth}
\begin{tabular}{lccccc}
\toprule
모델 & fold 0 & fold 1 & fold 2 & fold 3 & fold 4 \\
\midrule
베이스라인 nano & 93 / 73 & 91 / 71 & 89 / 69 & 100 / 86 & 98 / 78 \\
small & 100 / 97 & 91 / 71 & 100 / 89 & 100 / 84 & 100 / 100 \\
제안 (BRM nano) & 84 / 64 & 100 / 90 & 90 / 70 & 81 / 61 & 98 / 78 \\
\bottomrule
\end{tabular}
\end{adjustbox}
\end{table}


\section{교차 데이터셋 시드별 결과}
% 자동 생성 파일 — scripts/make_thesis.py 가 덮어쓴다. 직접 고치지 말 것.
\begin{table}[htbp]
\centering
\footnotesize
\caption[LODO 시드별 상세 결과 (베이스라인 / 제안).]{LODO 시드별 상세 결과 (베이스라인 / 제안). 체크포인트 epoch은 검증 fitness가 최대인 epoch, 총 epoch은 조기 종료까지 학습한 epoch 수이다.}
\label{tab:lodo_seed}
\begin{adjustbox}{max width=\linewidth}
\begin{tabular}{llccccc}
\toprule
방향 & 시드 & Mask mAP$_{50:95}$ & Dice & HD95 (\%) & 체크포인트 epoch & 총 epoch \\
\midrule
\multirow{3}{*}{TN3K$\rightarrow$DDTI} & 42 & 0.4501 / 0.4378 & 0.6954 / 0.5771 & 17.89 / 34.15 & 39 / 25 & 59 / 45 \\
 & 43 & 0.4580 / 0.4304 & 0.6809 / 0.6771 & 20.13 / 19.53 & 36 / 22 & 56 / 42 \\
 & 44 & 0.4315 / 0.4413 & 0.6331 / 0.6525 & 26.10 / 23.76 & 21 / 34 & 41 / 54 \\
\midrule
\multirow{3}{*}{DDTI$\rightarrow$TN3K} & 42 & 0.3070 / 0.2990 & 0.6152 / 0.5782 & 22.55 / 24.77 & 46 / 39 & 66 / 59 \\
 & 43 & 0.3384 / 0.2991 & 0.6498 / 0.5997 & 18.64 / 21.27 & 71 / 46 & 91 / 66 \\
 & 44 & 0.2985 / 0.3064 & 0.5578 / 0.6059 & 28.33 / 23.46 & 46 / 48 & 66 / 68 \\
\bottomrule
\end{tabular}
\end{adjustbox}
\end{table}


\chapter{구현 세부와 시행착오}\label{app:impl}

본 절은 같은 실험을 재현하거나 확장하려는 연구자를 위해, 구현 과정에서 실제로 겪은 문제와 해결 방법을
기록한다. 대부분 오류 메시지 없이 조용히 실패하는 문제였다는 점에서 기록할 가치가 있다고 판단하였다.

\begin{enumerate}
  \item \textbf{검출 헤드 양자화 시 출력 붕괴.} ONNX Runtime 정적 양자화를 그래프 전체에 적용하면 오류 없이
    변환되지만 모든 클래스 점수가 0이 된다. 헤드(\texttt{/model.23/} 이하 119개 노드)를 양자화 대상에서
    제외하여 해결하였다(\secref{sec:quant}).
  \item \textbf{채널별 양자화의 opset 요구.} 채널별 QDQ 양자화는 opset 12 그래프에서 유효하지 않은 그래프를
    만들어 로드되지 않는다. opset 13으로 변환하여 해결하였다.
  \item \textbf{OpenVINO 백엔드의 CUDA 비활성화.} 한 프로세스에서 OpenVINO 모델을 불러오면 그 뒤로 같은
    프로세스의 CUDA가 비활성화되어 GPU 평가가 모두 실패한다. 정밀도 단위로 평가 순서를 바꾸고 OpenVINO를
    마지막에 두어 회피하였다.
  \item \textbf{의존성 자동 설치에 의한 프로세스 종료.} 설치된 패키지가 \texttt{onnxruntime-gpu}여서 Ultralytics가
    ONNX Runtime이 없다고 판단하고 실행 중에 자동 설치를 시도하다 프로세스가 종료되었다. 환경 변수
    \texttt{YOLO\_AUTOINSTALL=false}로 자동 설치를 끄고 해결하였다.
  \item \textbf{ONNX 모델의 과제 추론.} \texttt{.onnx} 파일을 불러올 때 과제 종류가 검출로 잘못 추론되어
    마스크가 생성되지 않았다(경계 지표가 모두 0). 과제를 분할로 명시하여 해결하였다.
  \item \textbf{OpenVINO 모델 폴더 이름.} Ultralytics는 폴더 이름이 \texttt{\_openvino\_model}로 끝나야
    OpenVINO 모델로 인식한다.
  \item \textbf{학습 중복 실행에 의한 결과 오염.} 앞의 \appref{app:repro}에서 설명한 문제이다. 오염된 실행은
    학습 곡선의 epoch 수가 최대치(100)를 넘는 것으로 식별하였다.
\end{enumerate}
